\section{NVIDIA、Nemotron 3.5 Super VLを発表}
NVIDIAは長時間稼働するマルチモーダルエージェント向けのオープンモデル「Nemotron 3.5 Super VL」を発表した。MoE構成で総パラメータ120B、実際に使用するパラメータは12Bと説明し、推論・ツール利用に加えて文章、画像、文書、動画の理解を統合する。公式投稿は、10,000件のコーディングタスクを最大4倍高速に処理しつつSWE-Bench Proの精度を維持したと主張する。関連するBolt系列やNVFP4最適化版にも言及されたが、比較条件やローカル実行要件の詳細は今後の確認事項となる。
\dailyxsource{NVIDIA AI (@NVIDIAAI)}{https://x.com/NVIDIAAI/status/2110749836450713896}

\section{Meta、Muse Spark Openのウェイトを公開}
Metaは「Muse Spark Open」のオープンウェイトを公開した。総816B／アクティブ42BのMoEで、100万トークンの文脈長に加え、テキスト、画像、動画、音声を扱うマルチモーダル機能を掲げる。公式情報では、コーディングやエージェント用途、API提供経路のほか、Apple M5 Ultra搭載Mac Studio 2台でのローカル実行にも触れている。ライセンスはMuse Model Licenseとされ、利用条件には確認が必要である。「米国製オープンウェイトで最も有能」との位置付けは開発元の主張として扱う。
\dailyxsource{AI at Meta (@AIatMeta)}{https://x.com/AIatMeta/status/2110779754274292212}

\section{Anthropic、METRへの評価アクセスを拡大}
Anthropicは独立したAI評価機関との協力方法を更新し、METRのチームに従業員相当のアクセスを提供すると発表した。モデル開発側が外部評価者に継続的な検証環境を用意する取り組みであり、安全性評価の透明性と独立性に関わる。公式投稿と同社の説明をGrokが確認したと報告しているが、これによって得られる具体的な評価結果やMETR側の検証報告は本号の収集範囲に含まれていない。運用手続きの公表と、評価成果の公開は分けて捉える必要がある。
\dailyxsource{Anthropic (@AnthropicAI)}{https://x.com/AnthropicAI/status/2110784662071398807}

\section{Simons Foundation、数学者支援を発表}
Simons Foundationは、AI時代における数学コミュニティの活動を支援するため、1,500万ドルの取り組みを発表した。AnthropicとOpenAIがパートナーとなり、資金や計算資源、モデルへのアクセスを提供するという。Grokの記録には、OpenAIによる500万ドルの拠出に触れる関連投稿もある。ワークショップや会議、研究者向けリソースの整備が想定される一方、各支援プログラムの詳細や実際の研究への影響は今後の課題である。AIによる数学研究の変化を、研究者側から支える動きとして注目される。
\dailyxsource{Anthropic (@AnthropicAI)}{https://x.com/AnthropicAI/status/2110846112152105136}

\section{OpenAI、Rosalind Academic Planを拡大}
OpenAIは学術研究向けの「Rosalind Academic Plan」を拡大し、GPT-Rosalindによる生物学や科学文献、データの推論支援に利用できるスポンサー付きアクセスを案内した。Grokが確認した説明では、米国の大学研究室を対象とし、対象データを既定ではモデルの訓練に利用しないという条件にも触れている。研究者が専門領域でAIを試す際の費用や運用上の障壁を下げる狙いが見える。ただし、適格要件や実際の研究成果については、今回の投稿だけでは判断できない。
\dailyxsource{OpenAI (@OpenAI)}{https://x.com/OpenAI/status/2110841898076418336}

\section{FLUX 3 Image Klein、画像生成・編集に対応}
Black Forest Labsは、5.6Bパラメータのオープンウェイト画像モデル「FLUX 3 Image Klein」を紹介した。ネイティブ4Kの画像生成に加え、バウンディングボックスを用いた操作やピクセル単位の編集に対応し、単一のコンシューマーGPUで動作すると主張している。NVIDIAによる関連投稿では、学習基盤やNVFP4変種も取り上げられた。ただしGrokの収集では公式モデルカードなど一次資料の直接確認は限定的であり、対応機能、実行要件や画質は開発元の主張として記録する。
\dailyxsource{Black Forest Labs (@bfl\_ai)}{https://x.com/bfl_ai/status/2110751201105256925}

\section{Grok Botの利用拡大と自律タスク}
Elon Musk氏はGrok Botの利用が急速に伸びていると投稿し、無料ティアがなくても価値を提供できていると主張した。X上では飛行機の予約や迷惑投稿の除去など、自律的に作業を進める使い方についても言及があった。成長率や利用規模を裏付ける公式統計は、Grokの収集では確認されていない。なお、主な根拠とされた本人投稿は\textbf{10月16日07:52 JST}で、今回の観測窓（同日07:00まで）の\textbf{終了後}に当たる。厳密には窓外の参考情報として扱い、窓内の出来事とは区別する。
\dailyxsource{Elon Musk (@elonmusk)、観測窓外の投稿}{https://x.com/elonmusk/status/2110866513955868924}

\section{Ling-3.1など、その他のオープンモデル言及}
観測窓内のまとめ投稿では、Ant Groupの「Ling-3.1-flash」のオープンソース化や、Kunlun Wanweiの「Matrix-Game 4.0」などが紹介された。どちらもGrokが一次資料を確認したわけではなく、確認できたのは単一のラウンドアップ投稿が中心である。各モデルの正式名称、公開形態、実際の配布先や性能について、現段階では断定を避ける。公式発表やリポジトリが確認されるまでは、コミュニティで言及された話題として位置付ける。
\dailyxsource{Aaron.Z (@soluunj)}{https://x.com/soluunj/status/2110567964818964912}
